\section{Ascent to an arbitrary base field}\label{app:ascent}\label{sec:ascent}

Proposition~\ref{prop:ascent}, applied in Appendix~\ref{app:assembly} with $E_1$ and $F_1$ as in Theorem~\ref{thm:laurent},
carries analyticity at no point from a spherically complete Laurent subfield $K_1$ to the completion $K'$ of the prescribed
field, since a lift after scalar extension descends to $K_1$. Lemma~\ref{lem:descent} then passes from $K'$ to the field itself.

\subsection{The Laurent subfield}
We use spherical completeness in the sense used in Theorem~\ref{thm:spherical}: every nonempty family of pairwise-meeting
closed balls has a common point. For the field case of the following lemma, see
Schneider~\cite[Lemmas~1.3 and~1.6]{Sch}.

\begin{lemma}[discrete distances]\label{lem:discdist}
Let $Y$ be a complete nonarchimedean space whose distances lie in $r^\Z\cup\{0\}$ for some $r\in(0,1)$. Then $Y$ is
spherically complete.
\end{lemma}

\begin{proof}
Write $\cD(a,\rho)=\{y\in Y:\mathrm d(y,a)\le\rho\}$. In a nonarchimedean space, every point of a ball is a center.
Consequently, whenever two closed balls meet, the one with the smaller radius lies in the other.
The discrete distance assumption also gives
\[
\cD(a,\rho)=\cD(a,\rho^\flat),\qquad
\rho^\flat:=\max\{s\in r^\Z\cup\{0\}:s\le\rho\}.
\]
We may therefore replace every radius by a value in $r^\Z\cup\{0\}$.

Let $(\cD(a_i,\rho_i))_{i\in I}$ be a nonempty family of pairwise-meeting closed balls. If its radii have a least
element $\rho_{i_0}$, then $\cD(a_{i_0},\rho_{i_0})$ lies in every ball of the family, so $a_{i_0}$ is a common point.
Otherwise, discreteness implies that all the radii are positive and that their infimum is zero. Choose
$i_1,i_2,\dots$ so that $\rho_{i_n}$ decreases strictly to zero. The corresponding balls are nested, and
\[
\mathrm d(a_{i_{n'}},a_{i_n})\le\rho_{i_n}\qquad(n'\ge n).
\]
Their centers form a Cauchy sequence. By completeness it has a limit $a$, and closedness puts $a$ in every ball
of the sequence. For any $i\in I$, choose $n$ with $\rho_{i_n}\le\rho_i$. The $n$\nobreakdash-th ball then lies in
$\cD(a_i,\rho_i)$, so $a$ belongs to every ball of the original family.
\end{proof}



\begin{lemma}[Laurent subfield]\label{lem:laurentsub}
Let $K'$ be a complete nontrivially normed field of characteristic $p$, and let $t\in K'$ with $0<\abs t=r<1$.
\begin{enumerate}
\item $K'$ is nonarchimedean, and $\abs c=1$ for every $c\in\F_p^\times$.
\item For $a_m,\dots,a_M\in\F_p$, not all zero, $\bigl\lvert\sum_{i=m}^Ma_it^i\bigr\rvert=r^{\min\{i:a_i\ne0\}}$.
\item Give $\F_p((X))$ the absolute value with $\abs X=r$. There is a continuous isometric ring homomorphism
$\mathrm{ev}_t\colon\F_p((X))\to K'$ such that
$\mathrm{ev}_t\bigl(\sum_{i\ge m}a_iX^i\bigr)=\sum_{i\ge m}a_it^i$. Its image
$K_1$ is a closed subfield of $K'$; it is the closure of $\F_p(t)$, and it is isometrically isomorphic to $\F_p((X))$.
\item $K_1$ is spherically complete.
\end{enumerate}
\end{lemma}

\begin{proof}
(1) Lemma~\ref{lem:ultra} makes $K'$ nonarchimedean. For $c\in\F_p^\times$, the identity $c^{p-1}=1$ gives $\abs c=1$.

(2) The nonzero terms $a_it^i$ have pairwise distinct absolute values $r^i$. In a nonarchimedean field,
$\abs{x+y}=\abs x$ whenever $\abs y<\abs x$. Repeatedly applying this identity shows that the sum has the absolute
value of its largest term, namely $r^{\min\{i:a_i\ne0\}}$.

(3) Evaluation on Laurent polynomials defines a ring homomorphism
$\mathrm{ev}_t\colon\F_p[X,X^{-1}]\to K'$, and (2) shows that
$\abs{\mathrm{ev}_t(g)}=r^{\ord g}=\abs g$. Laurent polynomials are dense in $\F_p((X))$, since truncations converge.
Completeness of $K'$ therefore gives a unique continuous extension to $\F_p((X))$. Passing to limits preserves the
isometry and the ring operations. It also gives the stated evaluation formula, because the partial sums of
$\sum_{i\ge m}a_iX^i$ map to the partial sums of $\sum_{i\ge m}a_it^i$.

The image is complete as the isometric image of a complete space, hence closed in $K'$. It is a subfield: for
$g\ne0$, the identity $\mathrm{ev}_t(g)\mathrm{ev}_t(g^{-1})=1$ provides the inverse. This subfield contains
$\F_p(t)$, while the Laurent polynomials $\F_p[t,t^{-1}]\subseteq\F_p(t)$ are dense in it. The image is consequently
exactly the closure of $\F_p(t)$.

(4) By (3), $K_1$ is complete and has value group $\abs{K_1^\times}=r^\Z$. It is nonarchimedean by (1), so its distances
lie in $r^\Z\cup\{0\}$. Lemma~\ref{lem:discdist} gives spherical completeness.
\end{proof}

\subsection{Ingleton's theorem and its consequences}
Throughout the rest of this section we fix fields $K_1\subseteq K'$ with the following properties.
\begin{enumerate}
\item[(H1)] $K_1$, with the restricted absolute value, is nontrivially normed, complete and spherically complete.
\item[(H2)] $K'$ is nonarchimedean.
\end{enumerate}
By multiplication, $K'$ is a nonarchimedean normed $K_1$\nobreakdash-space. In our application $K'=\Kh$ and $K_1$ is the field of
Lemma~\ref{lem:laurentsub}.

\begin{theorem}[Ingleton~\cite{In}; see also~\cite{vR}]\label{thm:ingleton}
Let $K_1$ be a spherically complete nonarchimedean normed field, $V$ a $K_1$\nobreakdash-vector space with a nonarchimedean norm,
$V_0\subseteq V$ a subspace, $C\ge0$, and $T_0\colon V_0\to K_1$ linear with $\abs{T_0v}\le C\norm v$ for $v\in V_0$. Then
$T_0$ extends to a linear map $T\colon V\to K_1$ with $\abs{Tx}\le C\norm x$ for all $x\in V$.
\end{theorem}

Schneider~\cite[Prop.~9.2]{Sch} gives the more general formulation with a dominating nonarchimedean seminorm.

We will use Theorem~\ref{thm:ingleton} on the nonarchimedean $K_1$\nobreakdash-space $K'$ in two ways: to obtain the scalar
retraction below, and to extend finite-dimensional coordinate functionals in the proof of positivity of the
projective tensor norm. Both applications have codomain $K_1$. The spaces $E_1$ and $F_1$ introduced later need not
be nonarchimedean.

\begin{corollary}[Ingleton projection]\label{cor:proj}
There is a $K_1$\nobreakdash-linear map $\varpi\colon K'\to K_1$ with $\varpi|_{K_1}=\id$ and $\abs{\varpi(\lambda)}\le\abs\lambda$ for all
$\lambda\in K'$.
\end{corollary}

\begin{proof}
Theorem~\ref{thm:ingleton} with $V=K'$, $V_0=K_1$, $T_0=\id$ and $C=1$; its hypotheses hold by (H1) and (H2).
\end{proof}

\begin{samepage}
\begin{lemma}[extension]\label{lem:ext}
Let $\mathbb K\in\{K_1,K'\}$ and let $C\ge0$.
\begin{enumerate}
\item Let $V_0$ be a dense $\mathbb K$\nobreakdash-subspace of a normed $\mathbb K$\nobreakdash-space $V$, let $Y$ be a
$\mathbb K$\nobreakdash-Banach space and $T_0\colon V_0\to Y$ a $\mathbb K$\nobreakdash-linear map with $\norm{T_0x}\le C\norm x$. Then $T_0$ extends
uniquely to a continuous map $T\colon V\to Y$, which is $\mathbb K$\nobreakdash-linear with $\norm{Tx}\le C\norm x$.
\item Let $V_1,\dots,V_k$ be normed $K'$\nobreakdash-spaces with dense subspaces $V_i^0$, $Y$ a $K'$\nobreakdash-Banach space, and
$M\colon\prod_iV_i^0\to Y$ a $K'$\nobreakdash-multilinear map with $\norm{M(d)}\le C\prod_i\norm{d_i}$. Then $M$ extends uniquely to a
continuous $K'$\nobreakdash-multilinear map $\overline M\colon\prod_iV_i\to Y$ with $\norm{\overline M}\le C$.
\end{enumerate}
\end{lemma}
\end{samepage}

\begin{proof}
(1) The bound makes $T_0$ Lipschitz, hence uniformly continuous. Density of $V_0$ and completeness of $Y$ give a
unique continuous extension. Linearity and the norm bound pass to limits by continuity of addition, scalar
multiplication and the norm.

(2) The case $k=0$ is immediate. For $k\ge1$, multilinearity gives the telescoping estimate
\[
\norm{M(d)-M(d')}\le C\rho^{k-1}\sum_i\norm{d_i-d'_i}
\]
whenever $\norm{d_i},\norm{d'_i}\le\rho$ for every $i$, with $\rho>0$. For any $v\in\prod_iV_i$, choose tuples
$d^{(n)}\in\prod_iV_i^0$ converging to $v$. Their coordinates are uniformly bounded, so the estimate makes
$M(d^{(n)})$ Cauchy in $Y$. Define $\overline M(v)$ to be its limit. The same estimate, applied to two approximating
sequences, proves that this definition is independent of the choices and is continuous on bounded sets.
Multilinearity and the product norm bound pass to limits. Finally, density of $\prod_iV_i^0$ proves uniqueness.
\end{proof}

\subsection{Projective base change}
The next construction extends scalars while retaining the original norm and a contractive retraction whenever
the original space is complete. The scalar retraction from Corollary~\ref{cor:proj} is the essential ingredient;
the construction applies to arbitrary normed spaces.

For a normed $K_1$\nobreakdash-space $V$ let $V_{K'}:=V\otimes_{K_1}K'$, a $K'$\nobreakdash-vector space through the second factor, and put
\[
\norm u_\pi:=\inf\Bigl\{\sum_j\norm{x_j}\,\abs{\lambda_j}\ :\ u=\sum_jx_j\otimes\lambda_j\Bigr\},
\]
the infimum over finite decompositions, with an ordinary sum. For a bounded $K_1$\nobreakdash-linear $\psi\colon K'\to K_1$ let
$T_\psi\colon V_{K'}\to V$ be the $K_1$\nobreakdash-linear map with $T_\psi(x\otimes\lambda)=\psi(\lambda)x$; it exists because
$(x,\lambda)\mapsto\psi(\lambda)x$ is $K_1$\nobreakdash-bilinear. Let $\iota_V\colon V\to V_{K'}$, $x\mapsto x\otimes1$; it is
$K_1$\nobreakdash-linear, since $cx\otimes1=x\otimes c=c\cdot(x\otimes1)$ for $c\in K_1$.

\begin{lemma}[projective base change]\label{lem:bc}
Let $V$ be a normed $K_1$\nobreakdash-space; no nonarchimedean hypothesis is made on $V$.
\begin{enumerate}
\item $\norm{T_\psi u}\le\norm\psi\,\norm u_\pi$ for every bounded $K_1$\nobreakdash-linear $\psi\colon K'\to K_1$.
\item $\norm\cdot_\pi$ is a $K'$\nobreakdash-norm on $V_{K'}$.
\item $\iota_V$ is isometric.
\item Let $V\widehat\otimes_\pi K'$ be the completion of $(V_{K'},\norm\cdot_\pi)$. If $V$ is complete, then $T_\varpi$, with
$\varpi$ as in Corollary~\ref{cor:proj}, extends to a $K_1$\nobreakdash-linear map $\Pi_V\colon V\widehat\otimes_\pi K'\to V$ with
$\norm{\Pi_V}\le1$ and $\Pi_V\circ\iota_V=\id_V$.
\end{enumerate}
\end{lemma}

\begin{proof}
(1) For every decomposition $u=\sum_jx_j\otimes\lambda_j$,
\[
\norm{T_\psi u}=\Bigl\lVert\sum_j\psi(\lambda_j)x_j\Bigr\rVert
\le\sum_j\norm\psi\abs{\lambda_j}\norm{x_j}.
\]
Taking the infimum over decompositions proves the contraction estimate.

(2) Concatenating decompositions gives the triangle inequality. Multiplying their second factors by
$\mu\in K'$ gives $\norm{\mu u}_\pi\le\abs\mu\norm u_\pi$; for $\mu\ne0$, applying the same inequality with
$\mu^{-1}$ gives equality. The zero-scalar case is immediate.

It remains to prove positivity. Let $u\ne0$. Choose a finite tensor representation and rewrite its second
factors in a basis of their $K_1$\nobreakdash-span. This gives
\[
u=\sum_{j=1}^nx_j\otimes\lambda_j
\]
with $\lambda_1,\dots,\lambda_n$ linearly independent over $K_1$. Some $x_{j_0}$ is nonzero, since otherwise
$u=0$. On the finite-dimensional subspace $U:=\operatorname{span}_{K_1}\{\lambda_1,\dots,\lambda_n\}\subseteq K'$,
take the coordinate functional $\psi_0(\lambda_j)=\delta_{jj_0}$. It is bounded by some $C\ge0$ by
Lemma~\ref{lem:findim}, using completeness of $K_1$ in (H1). Apply Theorem~\ref{thm:ingleton} to the nonarchimedean
$K_1$\nobreakdash-space $K'$ to extend it to $\psi\colon K'\to K_1$ with $\norm\psi\le C$. Then
$T_\psi u=x_{j_0}$, so (1) gives
\[
0<\norm{x_{j_0}}\le C\norm u_\pi.
\]
Thus $\norm u_\pi>0$. This argument detects a nonzero vector coefficient by a functional on the \emph{scalar}
factor $K'$; it requires no separating continuous dual on $V$.

(3) The elementary tensor representation gives $\norm{x\otimes1}_\pi\le\norm x$. For the reverse inequality,
use $\varpi(1)=1$ and $\norm\varpi\le1$ in (1):
\[
\norm x=\norm{T_\varpi(x\otimes1)}\le\norm{x\otimes1}_\pi.
\]

(4) When $V$ is complete, the contraction $T_\varpi$ on the dense tensor subspace extends by
Lemma~\ref{lem:ext}(1), with $\mathbb K=K_1$, to a $K_1$\nobreakdash-linear contraction $\Pi_V$ with values in $V$.
On the embedded copy of $V$ it satisfies
$\Pi_V\iota_Vx=T_\varpi(x\otimes1)=\varpi(1)x=x$.
\end{proof}

We now fix the spaces to which the descent argument will be applied. Let $E_1$ be a normed $K_1$\nobreakdash-space
and $F_1$ a $K_1$\nobreakdash-Banach space, and put
\[
E:=E_1\widehat\otimes_\pi K',\qquad F:=F_1\widehat\otimes_\pi K',
\]
the completions of $((E_1)_{K'},\norm\cdot_\pi)$ and $((F_1)_{K'},\norm\cdot_\pi)$. They are $K'$\nobreakdash-Banach spaces;
$\iota_E\colon E_1\to E$ and $\iota_F\colon F_1\to F$ are $K_1$\nobreakdash-linear isometries, and $\Pi_F\colon F\to F_1$ is $K_1$\nobreakdash-linear
with $\norm{\Pi_F}\le1$ and $\Pi_F\circ\iota_F=\id_{F_1}$ (Lemma~\ref{lem:bc}).

\begin{lemma}[operators]\label{lem:ops}
For $f\in\cL_{K_1}(E_1,E_1)$ the map $f\otimes\id$ on $(E_1)_{K'}$ is $K'$\nobreakdash-linear with
$\norm{(f\otimes\id)u}_\pi\le\norm f\norm u_\pi$. Its extension $f_{K'}\in\cL_{K'}(E,E)$ satisfies $\norm{f_{K'}}\le\norm f$
and $f_{K'}\circ\iota_E=\iota_E\circ f$, and $f\mapsto f_{K'}$ is $K_1$\nobreakdash-linear.
\end{lemma}

\begin{proof}
For any decomposition $u=\sum_jx_j\otimes\lambda_j$, the image under $f\otimes\id$ has the decomposition
$\sum_jf(x_j)\otimes\lambda_j$, whose cost is at most
$\norm f\sum_j\norm{x_j}\abs{\lambda_j}$. Taking the infimum proves the norm bound on the algebraic tensor product.
Since $E$ is complete, Lemma~\ref{lem:ext}(1), with $\mathbb K=K'$, gives the continuous extension $f_{K'}$ with the
same bound. Its action on $x\otimes1$ proves the embedding identity. Additivity and $K_1$\nobreakdash-homogeneity in
$f$ hold on elementary tensors, hence on their span, and then on $E$ by continuity.
\end{proof}

\begin{lemma}[forms]\label{lem:forms-bc}
For $m\in\Alt^k_{K_1}(E_1;F_1)$ there is $m_{K'}\in\Alt^k_{K'}(E;F)$ with $\norm{m_{K'}}\le\norm m$ and
\[
m_{K'}(\iota_Ex_1,\dots,\iota_Ex_k)=\iota_F\bigl(m(x_1,\dots,x_k)\bigr)\qquad(x_i\in E_1),
\]
and $m\mapsto m_{K'}$ is $K_1$\nobreakdash-linear.
\end{lemma}

\begin{proof}
We first define the form on algebraic tensors, then establish a bound independent of their decompositions, and
finally pass to the completions. Fix a $K_1$\nobreakdash-basis $(\lambda_\beta)_{\beta\in\mathfrak B}$ of $K'$.
Every $u\in(E_1)_{K'}$ has a unique finite expansion
$u=\sum_\beta u^\beta\otimes\lambda_\beta$, with $u^\beta\in E_1$, and each coordinate map $u\mapsto u^\beta$ is
$K_1$\nobreakdash-linear. Define
\[
m^{\mathrm{alg}}(u_1,\dots,u_k):=
\sum_{\beta_1,\dots,\beta_k}
m\bigl(u_1^{\beta_1},\dots,u_k^{\beta_k}\bigr)\otimes
\lambda_{\beta_1}\cdots\lambda_{\beta_k}\in(F_1)_{K'}.
\]
The sum is finite, since each input has finitely many nonzero coordinates.

\textup{(i)} Linearity of the coordinate maps and multilinearity of $m$ make $m^{\mathrm{alg}}$
$K_1$\nobreakdash-multilinear.

\textup{(ii)} On elementary tensors, the definition gives
\[
m^{\mathrm{alg}}(x_1\otimes\mu_1,\dots,x_k\otimes\mu_k)
=m(x_1,\dots,x_k)\otimes\mu_1\cdots\mu_k
\qquad(x_i\in E_1,\ \mu_i\in K').
\]
Indeed, writing $\mu_i=\sum_\beta c_{i\beta}\lambda_\beta$ with $c_{i\beta}\in K_1$ gives
$x_i\otimes\mu_i=\sum_\beta(c_{i\beta}x_i)\otimes\lambda_\beta$. Substitution in the defining sum yields
\[
\sum_{\beta_1,\dots,\beta_k}
\Bigl(\prod_ic_{i\beta_i}\Bigr)m(x)\otimes\prod_i\lambda_{\beta_i}
=m(x)\otimes\prod_i\Bigl(\sum_\beta c_{i\beta}\lambda_\beta\Bigr),
\]
which is the asserted formula.

\textup{(iii)} The map $m^{\mathrm{alg}}$ is $K'$\nobreakdash-multilinear. By additivity in each slot, it suffices
to check $K'$\nobreakdash-homogeneity when every argument is an elementary tensor; there it follows from
\textup{(ii)}. The same formula determines the map on all tuples of algebraic tensors, so the initial basis
choice does not affect it.

\textup{(iv)} For arbitrary decompositions $u_i=\sum_jx_{ij}\otimes\mu_{ij}$, multilinearity and \textup{(ii)} give
\[
m^{\mathrm{alg}}(u)=
\sum_{j_1,\dots,j_k}m(x_{1j_1},\dots,x_{kj_k})\otimes\prod_i\mu_{ij_i}.
\]
The cost of this tensor decomposition is at most
$\norm m\prod_i\bigl(\sum_j\norm{x_{ij}}\abs{\mu_{ij}}\bigr)$.
Taking the infimum separately over the decompositions in each slot proves
\[
\norm{m^{\mathrm{alg}}(u_1,\dots,u_k)}_\pi\le\norm m\prod_i\norm{u_i}_\pi.
\]

\textup{(v)} To prove strong alternation, suppose two distinct slots $a,b$ both contain
$u=\sum_jx_j\otimes\mu_j$. Use this \emph{same} decomposition in both slots and expand all arguments into
elementary tensors. Terms using the same index $j$ in slots $a$ and $b$ vanish because $m$ has a repeated
argument. For $j\ne j'$, the term using indices $(j,j')$ in these slots pairs with the term using $(j',j)$.
Their scalar factors agree, while the values of $m$ are negatives of one another, so the pair sums to zero.
Thus $m^{\mathrm{alg}}$ vanishes whenever any two arguments coincide. Pair cancellation remains valid in
characteristic $2$; merely checking $m^{\mathrm{alg}}(u,\dots,u)=0$ would not establish strong alternation for
$k\ge3$.

\textup{(vi)} Regard $m^{\mathrm{alg}}$ as taking values in the completion $F$.
By \textup{(iii)}, \textup{(iv)} and Lemma~\ref{lem:ext}(2), it extends to a continuous
$K'$\nobreakdash-multilinear map $m_{K'}\colon E^k\to F$ with $\norm{m_{K'}}\le\norm m$.
Strong alternation also passes to the limit: if $v_a=v_b$, approximate this common vector by one sequence in
$(E_1)_{K'}$, and approximate the other entries separately. Every approximating tuple has equal entries in
slots $a,b$, so \textup{(v)} and continuity give $m_{K'}(v)=0$.

\textup{(vii)} Taking every scalar factor equal to one in \textup{(ii)} gives
\[
m_{K'}(\iota_Ex_1,\dots,\iota_Ex_k)
=m^{\mathrm{alg}}(x_1\otimes1,\dots,x_k\otimes1)
=m(x)\otimes1=\iota_F(m(x)).
\]
The defining formula is $K_1$\nobreakdash-linear in $m$. These linear identities pass to the continuous
extensions by uniqueness, proving $K_1$\nobreakdash-linearity of $m\mapsto m_{K'}$.
\end{proof}

\begin{proposition}[descent of a lift]\label{prop:ascent}
Assume \textup{(H1)} and \textup{(H2)}. Let $E_1$ be a normed $K_1$\nobreakdash-space, $F_1$ a $K_1$\nobreakdash-Banach space, and $E,F$ as
above. If $A^{k,K'}_{E,E;F}$ admits a bounded $k$\nobreakdash-linear lift $P$ over $K'$, then $A^{k,K_1}_{E_1,E_1;F_1}$ admits a
bounded $k$\nobreakdash-linear lift $P_1$ over $K_1$ with $\norm{P_1}\le\norm P$. Equivalently, by Proposition~\ref{prop:homogeneous-lift} over $K_1$ and over
$K'$: if $A^{k,K_1}_{E_1,E_1;F_1}$ is analytic at no point, then $A^{k,K'}_{E,E;F}$ is analytic at no point either.
\end{proposition}

\begin{proof}
Extend the operator and form inputs to $K'$, apply the given lift, and retract the value to $F_1$. Explicitly,
for $f_1,\dots,f_k\in\cL_{K_1}(E_1,E_1)$, $m\in\Alt^k_{K_1}(E_1;F_1)$ and $x\in E_1^k$, set
\[
P_1(f_1,\dots,f_k)(m)(x_1,\dots,x_k):=
\Pi_F\Bigl(P\bigl((f_1)_{K'},\dots,(f_k)_{K'}\bigr)(m_{K'})(\iota_Ex_1,\dots,\iota_Ex_k)\Bigr).
\]
We check the norm bound and then the diagonal identity.

For fixed $f$ and $m$, the map $P((f_i)_{K'})(m_{K'})$ is a continuous alternating $K'$\nobreakdash-multilinear form
on $E$. Precomposing each slot with the $K_1$\nobreakdash-linear isometry $\iota_E$ and postcomposing with
$\Pi_F$ gives a $K_1$\nobreakdash-multilinear map on $E_1$, still vanishing whenever two arguments coincide.
Lemmas~\ref{lem:bc}--\ref{lem:forms-bc} give
\[
\norm{P_1(f)(m)}
\le\norm{\Pi_F}\norm P\prod_i\norm{(f_i)_{K'}}\norm{m_{K'}}
\le\norm P\prod_i\norm{f_i}\norm m.
\]
In particular, $P_1(f)(m)\in\Alt^k_{K_1}(E_1;F_1)$.
The form extension is $K_1$\nobreakdash-linear in $m$, the operator extension is $K_1$\nobreakdash-linear in each
$f_i$, and $P$ is $K'$\nobreakdash-multilinear with values in $K'$\nobreakdash-linear operators.
Consequently $P_1$ is $K_1$\nobreakdash-linear in $m$ and $K_1$\nobreakdash-multilinear in $(f_1,\dots,f_k)$.
The displayed estimate proves boundedness at both levels and yields $\norm{P_1}\le\norm P$.

On the diagonal, use $P(g,\dots,g)=A^{k,K'}(g)$ with $g=f_{K'}$, followed by the embedding identities in
Lemmas~\ref{lem:ops} and~\ref{lem:forms-bc} and the retraction identity $\Pi_F\iota_F=\id$:
\begin{align*}
P_1(f,\dots,f)(m)(x)
&=\Pi_F\bigl(m_{K'}(f_{K'}\iota_Ex_1,\dots,f_{K'}\iota_Ex_k)\bigr)\\
&=\Pi_F\bigl(m_{K'}(\iota_Efx_1,\dots,\iota_Efx_k)\bigr)\\
&=\Pi_F\bigl(\iota_F(m(fx_1,\dots,fx_k))\bigr)\\
&=m(fx_1,\dots,fx_k)=A^{k,K_1}(f)(m)(x).\qedhere
\end{align*}
\end{proof}

\begin{remark}\label{rem:ascent}
Proposition~\ref{prop:ascent} is characteristic-free and uses no hypothesis on $k$. It transports negative results
upward only, from $K_1$ to $K'$. Ingleton's theorem is used only on $K'$, with no nonarchimedean norm hypothesis on $F_1$. The completeness of $F_1$ is used essentially: it is what makes the retraction $\Pi_F$ land in $F_1$.
\end{remark}

\subsection{Incomplete base fields}
The preceding construction takes place over the completion $\Kh$ of the prescribed field. To return to $K$,
we restrict scalars on the resulting Banach spaces. Density of $K$ in $\Kh$ forces every continuous
$K$\nobreakdash-multilinear map between these spaces to be $\Kh$\nobreakdash-multilinear, so a lift cannot appear under
this restriction.

\begin{lemma}\label{lem:descent}
Let $K$ be a nontrivially normed field with completion $\Kh$, and let $E,E',F$ be $\Kh$\nobreakdash-Banach spaces, regarded as
$K$\nobreakdash-Banach spaces by restriction of scalars (with the same norms). Then $\cL_K(E,E')=\cL_{\Kh}(E,E')$,
$\Alt^k_K(E;F)=\Alt^k_{\Kh}(E;F)$ and $\Alt^k_K(E';F)=\Alt^k_{\Kh}(E';F)$ with the same norms, $A^{k,K}_{E,E';F}=A^{k,\Kh}_{E,E';F}$,
and every bounded $k$\nobreakdash-linear lift of $A^{k,K}_{E,E';F}$ over $K$ is a bounded $k$\nobreakdash-linear lift of $A^{k,\Kh}_{E,E';F}$ over
$\Kh$.
\end{lemma}

\begin{proof}
Let $g$ be a continuous $K$\nobreakdash-linear map between normed spaces that already carry compatible
$\Kh$\nobreakdash-scalar actions. For $\lambda\in\Kh$, choose $\lambda_n\in K$ with $\lambda_n\to\lambda$.
Since $\norm{\lambda_nx-\lambda x}=\abs{\lambda_n-\lambda}\norm x$, continuity gives
\[
g(\lambda x)=\lim_n g(\lambda_nx)=\lim_n\lambda_ng(x)=\lambda g(x).
\]
Thus $g$ is $\Kh$\nobreakdash-linear. Applying the same argument with all but one slot fixed shows that every
continuous $K$\nobreakdash-multilinear map is $\Kh$\nobreakdash-multilinear. Alternation is the same equal-input
vanishing condition over either field.

The spaces of bounded maps therefore coincide as sets. Their operator norms agree as well, since the norms are
defined by the same pointwise bounds on the same vectors. They are $\Kh$\nobreakdash-Banach spaces, and
$A^{k,K}_{E,E';F}$ and $A^{k,\Kh}_{E,E';F}$ agree by their defining precomposition formula.
Finally, apply the slotwise density argument once more to a bounded $k$\nobreakdash-linear lift $P$ between these
map spaces. It becomes $\Kh$\nobreakdash-multilinear with the same norm and the same diagonal identity, and hence
is a lift over $\Kh$.
\end{proof}

\needspace{8\baselineskip}

